% Section 6: Limits and explicit non-claims

This section is placed before reproducibility deliberately. The shortest
honest summary of this artifact is the list of things it does not do.

\subsection{This artifact does NOT contain a new proof}

Theorem~1.1 of \cite{OpenAI2026Proof,Sawin2026Remarks} is proved by
those authors. The present paper does not re-derive Lemma~2.1,
Lemma~2.2, Proposition~2.3, or any other intermediate result. Where the
source papers say ``by the Golod--Shafarevich theorem'' or ``by the
Hajir--Maire tower-cutting argument'', we likewise take those classical
results as given.

\subsection{This artifact does NOT construct the infinite tower}

The existence of $L_j$ with $[L_j : \Q] \to \infty$ and bounded root
discriminant is supplied by Golod--Shafarevich
\cite{GolodShafarevich1964} together with the Hajir--Maire splitting
technique \cite{HajirMaire2001} adapted to pro-$2$ in
\cite{Sawin2026Remarks}. This is a non-constructive existence theorem.
No finite computation can exhibit the entire infinite tower. Phase~3
of the artifact stops at the boundary of computable
admissibility.

\subsection{This artifact has NOT been peer-reviewed}

No external mathematician has reviewed this code base. Our confidence
in the artifact rests on three independent signals:

\begin{enumerate}[leftmargin=*]
  \item the 59-test unit suite passes on CI (Ubuntu + Windows
        $\times$ Python 3.11 / 3.12);
  \item the central numerical reproduction matches the published
        $\approx 6.24 \cdot 10^{-38}$ to $0.014\%$ relative error
        (\S\ref{sec:eq22});
  \item the source code is small (\(\approx\) 1\,000 LOC across four
        modules) and intended to be readable by anyone who has read
        \cite{Sawin2026Remarks} \S2.
\end{enumerate}

None of these substitutes for mathematician review. The artifact is
offered as input to such review, not as a substitute.

\subsection{Reported $0.014$ and $0.0318$ are NOT in either formal PDF}

Secondary coverage (e.g.\ \cite{KalaiBlog2026}) reports numerical
values $\delta \approx 0.014$ and $\delta \approx 0.0318$ in the
informal contexts $n^{1.014}$ and $n^{1.0318}$. We verified by direct
text extraction with \texttt{pypdf} that neither of these literal
numerical values appears in either
\cite{OpenAI2026Proof} or \cite{Sawin2026Remarks}. The only rigorous
numerical lower bound that appears verbatim in either formal PDF is
$\delta_{\text{excess}} \approx 6.24 \cdot 10^{-38}$ from eq.~(2.2) of
\cite{Sawin2026Remarks}.

The artifact's \texttt{SOURCE\_CLAIMS} registry encodes this
three-way provenance explicitly:

\begin{itemize}[leftmargin=*]
  \item \texttt{openai\_index\_reported} ($\delta = 0.014$,
        \texttt{verified\_numerical\_in\_formal\_pdf = False}),
  \item \texttt{sawin\_rigorous\_lower\_bound}
        ($\delta = 6.24 \cdot 10^{-38}$,
        \texttt{verified\_numerical\_in\_formal\_pdf = True}),
  \item \texttt{sawin\_refinement\_reported} ($\delta = 0.0318$,
        \texttt{verified\_numerical\_in\_formal\_pdf = False}).
\end{itemize}

\subsection{This artifact does NOT improve the published bound}

We made no attempt to optimise the choice of $T$, $S$, $L_T$ or the
inner parameter $k$ to obtain a larger $\delta_{\text{excess}}$. The
explicit choice of \cite[\S2.1]{Sawin2026Remarks} is what we encode and
reproduce; the optimisation problem of finding $(T, S, k)$ with larger
admissible $\delta$ remains open and is out of scope.
